\documentclass[border=4pt]{standalone}
\usepackage{amsmath,amssymb,mathtools,xcolor,varwidth}
\definecolor{ellblue}{HTML}{1E6FE0}
\definecolor{radgreen}{HTML}{00A35A}
\definecolor{equalpurple}{HTML}{8E24AA}
\begin{document}
\begin{varwidth}{28em}
\large\bfseries
Draw $HI$ parallel to the tangent at $P$, meeting the diameter $EC$ produced.      The great focal property of the \textcolor{ellblue}{ellipse} is that $PS+PH=2AC$: the sum      of the two \textcolor{radgreen}{focal radii} is constant and equals the major axis. From      this Newton deduces the pivotal equality \textcolor{equalpurple}{$EP = AC$}, the semi-major      axis itself. This single relation binds the ellipse's focal geometry to the      measurable lengths the force calculation will demand.
\end{varwidth}
\end{document}
